% ═══════════════════════════════════════════════════════════════
% SPANISCH ARBEITSBLATT - Sequenz 07, Block b: Números 21-100
% ═══════════════════════════════════════════════════════════════
% Fachfarbe: #c41e3a (Karminrot)
% SW-Modus: swmodusfalse (Farbe)
% ═══════════════════════════════════════════════════════════════

\documentclass[11pt, a4paper]{article}

% ═══════════════════════════════════════════════════════════════
% SW-MODUS TOGGLE
% ═══════════════════════════════════════════════════════════════
\newif\ifswmodus
\swmodusfalse    % Standard: Farbe

% ═══════════════════════════════════════════════════════════════
% PAKETE
% ═══════════════════════════════════════════════════════════════

\usepackage[utf8]{inputenc}
\usepackage[T1]{fontenc}
\usepackage[ngerman]{babel}
\usepackage{helvet}
\renewcommand{\familydefault}{\sfdefault}

\usepackage[margin=2cm]{geometry}
\usepackage{enumitem}
\usepackage{tabularx}
\usepackage{booktabs}
\usepackage{multirow}

\usepackage{xcolor}
\usepackage{framed}
\usepackage{tcolorbox}
\tcbuselibrary{skins}

\usepackage{fancyhdr}
\usepackage{lastpage}
\usepackage{needspace}

% ═══════════════════════════════════════════════════════════════
% FARBEN
% ═══════════════════════════════════════════════════════════════

\definecolor{spanisch-color}{HTML}{c41e3a}
\definecolor{grau-color}{HTML}{888888}
\definecolor{hellgrau-color}{HTML}{f5f5f5}
\definecolor{orange-color}{HTML}{b35c00}

\definecolor{sw-dark}{HTML}{000000}
\definecolor{sw-gray}{HTML}{666666}
\definecolor{sw-white}{HTML}{ffffff}

\ifswmodus
    \colorlet{spanisch}{sw-dark}
    \colorlet{grau}{sw-gray}
    \colorlet{hellgrau}{sw-white}
    \colorlet{orange}{sw-dark}
\else
    \colorlet{spanisch}{spanisch-color}
    \colorlet{grau}{grau-color}
    \colorlet{hellgrau}{hellgrau-color}
    \colorlet{orange}{orange-color}
\fi

\newcommand{\fachfarbe}{spanisch}

% ═══════════════════════════════════════════════════════════════
% VARIABLEN
% ═══════════════════════════════════════════════════════════════

\newcommand{\thema}{Números 21-100 y precios}
\newcommand{\sequenz}{07}
\newcommand{\block}{b}
\newcommand{\fach}{Spanisch}
\newcommand{\klasse}{7}
\newcommand{\maxpunkte}{20}
\newcommand{\datum}{\today}

% ═══════════════════════════════════════════════════════════════
% KOPF- UND FUSSZEILE
% ═══════════════════════════════════════════════════════════════

\pagestyle{fancy}
\fancyhf{}
\renewcommand{\headrulewidth}{0.4pt}
\renewcommand{\footrulewidth}{0.4pt}

\lhead{\fach{} -- Klasse \klasse}
\chead{AB-\sequenz\block}
\rhead{\datum}

\lfoot{}
\cfoot{}
\rfoot{Seite \thepage\ von \pageref{LastPage}}

% ═══════════════════════════════════════════════════════════════
% BEFEHLE
% ═══════════════════════════════════════════════════════════════

\newcommand{\punktebox}[1]{\hfill\fbox{\small \_/#1}}

\newcommand{\luecke}[1]{\rule{#1}{0.4pt}}

\newtcolorbox{merkkasten}[1][]{
    colback=hellgrau,
    colframe=\fachfarbe,
    colbacktitle=white,
    coltitle=\fachfarbe,
    boxrule=1pt,
    arc=0pt,
    left=8pt, right=8pt, top=8pt, bottom=8pt,
    fonttitle=\bfseries,
    attach boxed title to top left={yshift=-2mm, xshift=5mm},
    boxed title style={boxrule=0pt, colback=white},
    title=#1
}

\newtcolorbox{vokabelkasten}{
    colback=white,
    colframe=\fachfarbe,
    boxrule=0.5pt,
    arc=0pt,
    left=5pt, right=5pt, top=5pt, bottom=5pt
}

\newtcolorbox{hinweis}{
    colback=white,
    colframe=orange,
    colbacktitle=white,
    coltitle=orange,
    boxrule=1pt,
    arc=0pt,
    left=5pt, right=5pt, top=5pt, bottom=5pt,
    fonttitle=\bfseries,
    attach boxed title to top left={yshift=-2mm, xshift=5mm},
    boxed title style={boxrule=0pt, colback=white},
    title=Hinweis
}

\newenvironment{aufgabe}[1]{%
    \needspace{5\baselineskip}%
    \nopagebreak[4]%
    \vspace{10pt}
    \noindent\textbf{\textcolor{\fachfarbe}{Aufgabe #1}}
    \nopagebreak[4]%
    \vspace{5pt}
    \nopagebreak[4]%
    \begin{list}{}{\leftmargin=0pt}
    \item[]
}{%
    \end{list}
}

\newlist{teilaufgaben}{enumerate}{2}
\setlist[teilaufgaben,1]{label=\alph*), leftmargin=2em}
\setlist[teilaufgaben,2]{label=\roman*), leftmargin=2em}

\newcommand{\es}[1]{\textcolor{\fachfarbe}{\textbf{#1}}}

% ═══════════════════════════════════════════════════════════════
% DOKUMENT
% ═══════════════════════════════════════════════════════════════

\begin{document}

% ═══════════════════════════════════════════════════════════════
% KOPFBEREICH
% ═══════════════════════════════════════════════════════════════

\begin{center}
    {\Large\bfseries\textcolor{\fachfarbe}{Arbeitsblatt: \thema}}
\end{center}

\vspace{5pt}

\noindent
\begin{tabularx}{\textwidth}{|X|X|X|}
\hline
\textbf{Name:} \luecke{4cm} &
\textbf{Klasse:} \klasse &
\textbf{Datum:} \luecke{2cm} \\
\hline
\end{tabularx}

\vspace{15pt}

% ═══════════════════════════════════════════════════════════════
% MERKKASTEN 1: Zahlen 21-29
% ═══════════════════════════════════════════════════════════════

\begin{merkkasten}[Merkkasten 1: Die Zahlen 21-29]
\textbf{Achtung:} Die Zahlen 21-29 werden als \textbf{ein Wort} geschrieben!

\vspace{0.5em}

\begin{tabular}{llll}
21 = \luecke{3cm} & 22 = \luecke{3cm} & 23 = \luecke{3cm} \\[0.8em]
24 = \luecke{3cm} & 25 = \luecke{3cm} & 26 = \luecke{3cm} \\[0.8em]
27 = \luecke{3cm} & 28 = \luecke{3cm} & 29 = \luecke{3cm} \\
\end{tabular}

\vspace{0.5em}

\textit{Tipp: Übertrage die Zahlen aus dem Lernpfad!}
\end{merkkasten}

\vspace{10pt}

% ═══════════════════════════════════════════════════════════════
% MERKKASTEN 2: Zehner
% ═══════════════════════════════════════════════════════════════

\begin{merkkasten}[Merkkasten 2: Die Zehner 30-100]
\begin{tabular}{llll}
30 = \luecke{3cm} & 40 = \luecke{3cm} & 50 = \luecke{3cm} & 60 = \luecke{3cm} \\[0.8em]
70 = \luecke{3cm} & 80 = \luecke{3cm} & 90 = \luecke{3cm} & 100 = \luecke{3cm} \\
\end{tabular}

\vspace{0.5em}

\textbf{Kombinationen:} Zehner + \es{y} + Einer

Beispiel: 35 = treinta \es{y} cinco
\end{merkkasten}

\vspace{10pt}

% ═══════════════════════════════════════════════════════════════
% AUFGABE 1: Zahlen zuordnen
% ═══════════════════════════════════════════════════════════════

\begin{aufgabe}{1}
Verbinde die Zahl mit dem spanischen Wort. \punktebox{4}
\end{aufgabe}

\begin{center}
\begin{tabular}{rcl}
25 & \luecke{4cm} & sesenta y tres \\[0.8em]
47 & \luecke{4cm} & veinticinco \\[0.8em]
63 & \luecke{4cm} & ochenta y uno \\[0.8em]
81 & \luecke{4cm} & cuarenta y siete \\
\end{tabular}
\end{center}

\vspace{15pt}

% ═══════════════════════════════════════════════════════════════
% AUFGABE 2: Zahlen schreiben
% ═══════════════════════════════════════════════════════════════

\begin{aufgabe}{2}
Schreibe die Zahlen auf Spanisch. \punktebox{4}
\end{aufgabe}

\noindent
a) 22 = \luecke{5cm} \\[0.8em]
b) 38 = \luecke{5cm} \\[0.8em]
c) 54 = \luecke{5cm} \\[0.8em]
d) 99 = \luecke{5cm} \\

\vspace{15pt}

% ═══════════════════════════════════════════════════════════════
% MERKKASTEN 3: Preise
% ═══════════════════════════════════════════════════════════════

\begin{merkkasten}[Merkkasten 3: Los precios (Die Preise)]
\begin{tabular}{ll}
\es{¿Cuánto cuesta?} & = \luecke{5cm} \\[0.8em]
\es{Cuesta X euros.} & = \luecke{5cm} \\[0.8em]
\es{Cuesta X euros y Y céntimos.} & = \luecke{5cm} \\
\end{tabular}
\end{merkkasten}

\vspace{10pt}

% ═══════════════════════════════════════════════════════════════
% AUFGABE 3: Preise lesen
% ═══════════════════════════════════════════════════════════════

\begin{aufgabe}{3}
Lies die Preise und schreibe sie als Zahl. \punktebox{4}
\end{aufgabe}

\noindent
a) ``Cuesta veintitrés euros.'' = \luecke{2cm} € \\[0.8em]
b) ``Cuesta cuarenta y cinco euros.'' = \luecke{2cm} € \\[0.8em]
c) ``Cuesta setenta y ocho euros y cincuenta céntimos.'' = \luecke{2cm} € \\[0.8em]
d) ``Cuesta noventa y nueve euros.'' = \luecke{2cm} € \\

\vspace{15pt}

% ═══════════════════════════════════════════════════════════════
% AUFGABE 4: Preise sagen
% ═══════════════════════════════════════════════════════════════

\begin{aufgabe}{4}
Schreibe die Preise auf Spanisch. \punktebox{4}
\end{aufgabe}

\noindent
a) 35 € = Cuesta \luecke{6cm} euros. \\[0.8em]
b) 67 € = Cuesta \luecke{6cm} euros. \\[0.8em]
c) 82,50 € = Cuesta \luecke{4cm} euros y \luecke{4cm} céntimos. \\[0.8em]
d) 100 € = Cuesta \luecke{6cm} euros. \\

\vspace{15pt}

% ═══════════════════════════════════════════════════════════════
% AUFGABE 5: Dialog vervollständigen
% ═══════════════════════════════════════════════════════════════

\begin{aufgabe}{5}
Vervollständige den Einkaufsdialog. \punktebox{4}
\end{aufgabe}

\begin{hinweis}
Verwende: ¿Cuánto cuesta? | Cuesta ... euros | Gracias | aquí tiene
\end{hinweis}

\vspace{0.5em}

\noindent
\textbf{Cliente:} ¡Buenos días! \luecke{5cm} la camiseta?

\textbf{Vendedor:} \luecke{8cm}.

\textbf{Cliente:} Vale. \luecke{5cm}.

\textbf{Vendedor:} Muy bien. ¡\luecke{3cm}!

\vspace{20pt}

% ═══════════════════════════════════════════════════════════════
% PUNKTEÜBERSICHT
% ═══════════════════════════════════════════════════════════════

\hrule
\vspace{10pt}

\begin{flushright}
\textbf{Punkte:} \luecke{1cm} / \maxpunkte
\end{flushright}

\end{document}
